\documentclass[11pt,letter,notitlepage]{article}
\usepackage[left=2cm, right=2cm, lines=45, top=0.8in, bottom=0.7in]{geometry}
\usepackage{fancyhdr}
\usepackage{fancybox}
\usepackage{graphicx}
\usepackage{pdfpages}
\usepackage{enumitem}
\usepackage{algorithm}
\usepackage{algorithmic}
\renewcommand{\headrulewidth}{1.5pt}
\renewcommand{\footrulewidth}{1.5pt}
\newcommand\Loadedframemethod{TikZ}
\usepackage[framemethod=\Loadedframemethod]{mdframed}
\usepackage{amssymb,amsmath}
\usepackage{amsthm}
\usepackage{thmtools}
\usepackage{bm}
\setlength{\topmargin}{0pt}
\setlength{\textheight}{9in}
\setlength{\headheight}{0pt}
\setlength{\oddsidemargin}{0.25in}
\setlength{\textwidth}{6in}
%%%%%%%%%%%%%%%%%%%%%%%%
%%%%%% Define math operator %%%%%
%%%%%%%%%%%%%%%%%%%%%%%%
\DeclareMathOperator*{\argmin}{\bf argmin}
\DeclareMathOperator*{\relint}{\bf relint\,}
\DeclareMathOperator*{\dom}{\bf dom\,}
\DeclareMathOperator*{\intp}{\bf int\,}
%%%%%%%%%%%%%%%%%%%%%%%%
\setlength{\topmargin}{0pt}
\setlength{\textheight}{9in}
\setlength{\headheight}{0pt}
\setlength{\oddsidemargin}{0.25in}
\setlength{\textwidth}{6in}
\pagestyle{fancy}
%%%%%%%%%%%%%%%%%%%%%%%%
%% Define the Exercise environment %%
%%%%%%%%%%%%%%%%%%%%%%%%
\mdtheorem[
topline=false,
rightline=false,
leftline=false,
bottomline=false,
leftmargin=-10,
rightmargin=-10
]{exercise}{\textbf{Exercise}}
%%%%%%%%%%%%%%%%%%%%%%%%
%% Define the Problem environment %%
%%%%%%%%%%%%%%%%%%%%%%%%
\mdtheorem[
topline=false,
rightline=false,
leftline=false,
bottomline=false,
leftmargin=-10,
rightmargin=-10
]{problem}{\textbf{Problem}}
%%%%%%%%%%%%%%%%%%%%%%%%
%% Define the Solution Environment %%
%%%%%%%%%%%%%%%%%%%%%%%%
\declaretheoremstyle
[
spaceabove=0pt,
spacebelow=0pt,
headfont=\normalfont\bfseries,
notefont=\mdseries,
notebraces={(}{)},
headpunct={:\;},
headindent={},
postheadspace={ },
bodyfont=\normalfont,
qed=$\blacksquare$,
preheadhook={\begin{mdframed}[style=myframedstyle]},
postfoothook=\end{mdframed},
]{mystyle}
\declaretheorem[style=mystyle,title=Solution]{solution}
\mdfdefinestyle{myframedstyle}{%
	topline=false,
	rightline=false,
	leftline=false,
	bottomline=false,
	skipabove=-6ex,
	leftmargin=-10,
	rightmargin=-10,
	backgroundcolor=black!5}
%%%%%%%%%%%%%%%%%%%%%%%%

\renewcommand{\eqref}[1]{Eq.~(\ref{#1})}
\newcommand{\figref}[1]{Fig.~\ref{#1}}
\newcommand{\proj}[2]{\textbf{P}_{#2} (#1)}
\newcommand{\lspan}[1]{\textbf{span}  (#1)  }
\newcommand{\rank}[1]{ \textbf{rank}  (#1)  }
\newcommand{\tr}{ \operatorname{tr}  }
\newcommand{\RNum}[1]{\uppercase\expandafter{\romannumeral #1\relax}}
\newcommand{\mb}[1]{\mathbf{#1}}
\DeclareMathOperator*{\cl}{\bf cl\,}
\DeclareMathOperator*{\bd}{\bf bd\,}
\DeclareMathOperator*{\conv}{\bf conv\,}
\DeclareMathOperator*{\epi}{\bf epi\,}
\theoremstyle{definition}	
\newtheorem{definition}{Definition}	

%%%%%%%%%%%%%%%%%%%%%%%%%%%%%%%%%%%%%%%%%%
%%           Homework info.             %%
\newcommand{\posted}{\text{Nov. 28th, 2025}}
\newcommand{\due}{\text{Dec. 16th, 2025}}
\newcommand{\hwno}{\text{3}}
%%%%%%%%%%%%%%%%%%%%%%%%%%%%%%%%%%%%%%%%%%

%%%%%%%%%%%%%%%%%%%%%%%%%%%%%%%%%%%%%
%%    Put your information here   %%
\newcommand{\name}{\text{Jinghao Chen}}
\newcommand{\id}{\text{PB23010359}}
%%%%%%%%%%%%%%%%%%%%%%%%%%%%%%%%%%%%

\chead{\textbf{Homework \hwno}}

\begin{document}
	\vspace*{-4\baselineskip}
	\thispagestyle{empty}
	
	\begin{center}
		{\bf\large Introduction to Machine Learning}\\
		{Fall 2025}\\
		University of Science and Technology of China
	\end{center}
	
	\noindent
	Lecturer: Zhihui Li, Xiaojun Chang
	\hfill
	Homework \hwno
	\\
	Posted: \posted
	\hfill
	Due: \due
    \\
    Name: \name
	\hfill
	ID: \id
	
	\noindent
	\rule{\textwidth}{2pt}
	
	\medskip
	
	\textbf{Notice, }to get the full credits, please present your solutions step by step.

    \begin{exercise}[Affine Sets and Projections]
        \begin{enumerate}[label=(\arabic*)]
        \item Let \(A\in\mathbb{R}^{m\times n}\), \(b\in\mathbb{R}^{m}\), and consider
        \[U=\{x\in\mathbb{R}^{n}:Ax=b\}.\]
        \begin{enumerate}[label=(\alph*)]
            \item Show that \(U\) is an affine set. Identify its direction space in terms of \(A\).
            \item Prove that for any \(x_{0}\in U\),
            \[U=x_{0}+\mathcal{N}(A),\]
            where \(\mathcal{N}(A)\) denotes the null space of \(A\).
        \end{enumerate}
        
        \item Suppose \(A\in\mathbb{R}^{m\times n}\) has full row rank and \(U=\{x:Ax=b\}\) is nonempty.
        \begin{enumerate}[label=(\alph*)]
            \item Consider the projection of a point \(y\in\mathbb{R}^{n}\) onto \(U\):
            \[p=\argmin_{x\in U}\|x-y\|^{2}_{2}.\]
            Use Lagrange multipliers to derive the formula
            \[p=y-A^{\top}(AA^{\top})^{-1}(Ay-b).\]
            \item In \(\mathbb{R}^{3}\), let
            \[U=\{x\in\mathbb{R}^{3}:x_{1}+2x_{2}-x_{3}=3\},\quad y=(1,0,0)^{\top}.\]
            Compute explicitly the projection \(p\) of \(y\) onto \(U\).
        \end{enumerate}
        \end{enumerate}
    \end{exercise}

    \newpage

    \begin{exercise}[Convex Sets, Images and Preimages]
        Let \(C\subseteq\mathbb{R}^{n}\) be a nonempty convex set and let \(A\in\mathbb{R}^{m\times n}\), \(a\in\mathbb{R}^{m}\).
        \begin{enumerate}[label=(\arabic*)]
        \item Show that the affine image of \(C\),
        \[D=\{y\in\mathbb{R}^{m}:y=Ax+a,\,x\in C\},\]
        is convex.
        
        \item Define the affine preimage of \(C\) by
        \[E=\{y\in\mathbb{R}^{m}:Ay+a\in C\}.\]
        Show that \(E\) is convex.
        
        \item (Epigraph and hypograph)
        \begin{enumerate}[label=(\alph*)]
            \item Let \(f:\mathbb{R}^{n}\to\mathbb{R}\) be a convex function. The epigraph of \(f\) is 
            \[\epi f=\{(x,t)\in\mathbb{R}^{n+1}:t\geq f(x)\}.\]
            Show that \(\epi f\) is a convex set.
            \item Let \(g:\mathbb{R}^{n}\to\mathbb{R}\) be concave. Its hypograph is 
            \[\operatorname{hypo}g=\{(x,t)\in\mathbb{R}^{n+1}:t\leq g(x)\}.\]
            Show that \(\operatorname{hypo}g\) is convex.
        \end{enumerate}
        
        \item (Interior and relative interior -- computation) You do not need to prove your answers, but you should explain them briefly.
        \begin{enumerate}[label=(\alph*)]
            \item Find \(\intp C\) and \(\relint C\) for 
            \[C=\{x\in\mathbb{R}^{3}:x_{1}^{2}+x_{2}^{2}<4,\ x_{3}\geq 0\}.\]
            \item Let 
            \[K=\{x\in\mathbb{R}^{n}:x_{i}\geq 0,\ \sum_{i=1}^{n}x_{i}=1\}\]
            be the standard probability simplex. Determine \(\intp K\) (as a subset of \(\mathbb{R}^{n}\)) and \(\relint K\).
        \end{enumerate}
        \end{enumerate}
    \end{exercise}

    \newpage

    \begin{exercise}[Cones and Farkas-type Results]
        Let \(A=(a_{1},\ldots,a_{n})\in\mathbb{R}^{m\times n}\) and consider the conic hull
        \[\operatorname{cone}(A)=\left\{\sum_{i=1}^{n}\alpha_{i}a_{i}:\alpha_{i}\geq 0\right\}\subseteq\mathbb{R}^{m}.\]
        
        \begin{enumerate}[label=(\arabic*)]
        \item Show that \(\operatorname{cone}(A)\) is a convex cone, i.e.,
        \[x,y\in\operatorname{cone}(A),\ \lambda,\mu\geq 0\ \Rightarrow\ \lambda x+\mu y\in\operatorname{cone}(A).\]
        
        \item Show that \(\operatorname{cone}(A)\) is closed.
        
        \item (Gordan's theorem) Let \(A\in\mathbb{R}^{m\times n}\). Prove that exactly one of the following two systems has a solution:
        \begin{itemize}
            \item There exists \(x\in\mathbb{R}^{n}\) such that \(Ax>0\) (componentwise).
            \item There exists \(y\in\mathbb{R}^{m}\), \(y\geq 0\), \(y\neq 0\), such that \(A^{\top}y=0\).
        \end{itemize}
        Hint: Consider appropriate cones and use separation theorems (you may assume standard separation results for closed convex cones).
        
        \item Show how Farkas' Lemma can be deduced from Gordan's theorem in part (3).
        \end{enumerate}
    \end{exercise}

    \newpage
    \bibliography{refs}
    \bibliographystyle{abbrv}
	
\end{document}